\chapter*{Glosario de Términos}
% La siguiente línea es para que el glosario aparezca en el Índice General
\addcontentsline{toc}{chapter}{Glosario de Términos}

A continuación, se definen los principales términos técnicos, acrónimos y conceptos administrativos utilizados a lo largo de este Trabajo de Fin de Grado para facilitar su comprensión:

\begin{description}
    \item[AENA] Aeropuertos Españoles y Navegación Aérea. Entidad pública empresarial encargada de la gestión de los aeropuertos de interés general en España.
    \item[AESA] Agencia Estatal de Seguridad Aérea. Organismo del Estado que vela por que se cumplan las normas de aviación civil.
    \item[API] (\textit{Application Programming Interface}). Conjunto de subrutinas, funciones y procedimientos que ofrece cierta biblioteca para ser utilizada por otro software como una capa de abstracción.
    \item[AVSAF] (\textit{Aviation Safety}). Certificación de seguridad operacional requerida para el personal que accede a la zona de movimiento de un aeropuerto.
    \item[AVSEC] (\textit{Aviation Security}). Certificación centrada en la protección de la aviación civil contra actos de interferencia ilícita.
    \item[PCAP] Pliego de Cláusulas Administrativas Particulares. Documento que rige los pactos y condiciones definidores de los derechos y obligaciones de las partes en un contrato público.
    \item[PPT] Pliego de Prescripciones Técnicas. Documento que contiene las especificaciones técnicas que han de regir la realización de la prestación objeto de una licitación.
    \item[RPA] (\textit{Robotic Process Automation}). Automatización robótica de procesos mediante el uso de software con inteligencia artificial y capacidades de aprendizaje automático para manejar tareas repetitivas de gran volumen.
    \item[Webhook] Callback HTTP definido por el usuario que se activa mediante un evento específico, utilizado para notificar a sistemas externos (como n8n) en tiempo real.
    \item[ERP] (\textit{Enterprise Resource Planning}). Sistema de planificación de recursos empresariales. Software monolítico que integra y automatiza las principales prácticas y procesos de negocio de una empresa (contabilidad, recursos humanos, inventario, etc.) en una única plataforma.
    \item[ONTSI] Observatorio Nacional de Tecnología y Sociedad de la Información. Órgano colegiado de carácter consultivo de la Administración General del Estado de España, fuente de los datos estadísticos sobre digitalización de las pymes.
    \item[Open Source] (Código abierto). Modelo de desarrollo de software basado en la colaboración abierta, donde el código fuente es accesible públicamente para que cualquiera pueda estudiarlo, modificarlo y distribuirlo.
    \item[Backend] (o Back-end). Capa de acceso a datos y lógica de negocio de una aplicación o sistema. Opera en el servidor, procesando la información y comunicándose con la base de datos, por lo que no es directamente visible para el usuario final.
    \item[Base de datos relacional] Tipo de base de datos que almacena y proporciona acceso a puntos de datos que están relacionados entre sí, organizando la información habitualmente en tablas compuestas por filas y columnas.
    \item[JSON] (\textit{JavaScript Object Notation}). Formato de texto ligero, estructurado y legible por humanos, utilizado como estándar para el intercambio de datos entre diferentes sistemas y lenguajes de programación.
    \item[Kanban] Marco de trabajo visual para la gestión ágil de proyectos. Utiliza tableros divididos en columnas (que representan fases del proceso) y tarjetas (que representan tareas) para optimizar el flujo de trabajo y evitar cuellos de botella.
    \item[Script] Archivo o fragmento de código que contiene una secuencia de comandos o instrucciones diseñadas para ser ejecutadas de forma automatizada por un programa interpretador, en lugar de ser compiladas.
    \item[Web Scraping] Técnica utilizada para extraer información y datos de sitios web de forma automatizada. Consiste en el uso de software que simula la navegación humana para leer y capturar el contenido del código HTML de las páginas.
\end{description}